%% LyX 2.4.4 created this file.  For more info, see https://www.lyx.org/.
%% Do not edit unless you really know what you are doing.
\documentclass{article}
\usepackage{amsmath}
\usepackage{amssymb}
\usepackage{fontspec}
\setlength{\parindent}{0bp}
\usepackage{color}
\usepackage{float}
\usepackage{units}
\usepackage{geometry}
\geometry{verbose,tmargin=2.7cm,bmargin=2cm,lmargin=2cm,rmargin=2cm,headheight=2cm,headsep=0cm,footskip=0cm}
\usepackage{setspace}
\usepackage{esint}
\usepackage{microtype}
\onehalfspacing
\usepackage[pdfusetitle,
 bookmarks=true,bookmarksnumbered=true,bookmarksopen=true,bookmarksopenlevel=2,
 breaklinks=true,pdfborder={0 0 0},pdfborderstyle={},backref=section,colorlinks=true]
 {hyperref}
\hypersetup{
 linkcolor=black, citecolor=black}

\makeatletter

%%%%%%%%%%%%%%%%%%%%%%%%%%%%%% LyX specific LaTeX commands.
\providecommand\textquotedblplain{%
  \bgroup\addfontfeatures{Mapping=}\char34\egroup}

\@ifundefined{date}{}{\date{}}
%%%%%%%%%%%%%%%%%%%%%%%%%%%%%% User specified LaTeX commands.
% Engine guard: use XeLaTeX/LuaLaTeX
\usepackage{iftex}
\ifPDFTeX
  \PackageError{template}{Please compile with XeLaTeX or LuaLaTeX}{}
\fi

% ======== ADVANCED HEBREW ACADEMIC TEMPLATE ========
% Packages (trimmed to avoid conflicts)
% \usepackage{autobreak}  % REMOVE: can conflict with amsmath; use \allowdisplaybreaks
% \usepackage[T1]{fontenc}  % REMOVE: not needed with fontspec
\usepackage{relsize}
\usepackage{xcolor}
\usepackage{amssymb}
\usepackage{titlesec}
\usepackage{polyglossia}
\usepackage{fancyhdr}
\usepackage{tcolorbox}
\usepackage{tikz}
\usepackage{pgfplots}
\usepackage{booktabs}
\usepackage{array}
\usepackage{multirow}
\usepackage{enumitem}
\usepackage{float}
\usepackage{caption}
\usepackage{subcaption}
\usepackage{geometry}
\usepackage{microtype}
\usepackage{fontspec}
\usepackage{physics}
\usepackage{siunitx}
\usepackage{mathtools}
\usepackage{empheq}
\usepackage{mdframed}

% Colors (deduplicated)
\definecolor{primaryblue}{RGB}{33,150,243}
\definecolor{secondaryblue}{RGB}{63,81,181}
\definecolor{accentgreen}{RGB}{76,175,80}
\definecolor{warningcolor}{RGB}{255,152,0}
\definecolor{errorcolor}{RGB}{244,67,54}
\definecolor{infocolor}{RGB}{3,169,244}
\definecolor{notecolor}{RGB}{255,193,7}
\definecolor{theoremcolor}{RGB}{156,39,176}
\definecolor{definitioncolor}{RGB}{255,87,34}
\definecolor{examplecolor}{RGB}{0,150,136}
\definecolor{lightgray}{RGB}{245,245,245}
\definecolor{darkgray}{RGB}{66,66,66}
% Add aliases for colors referenced in layout to avoid missing-color errors
\colorlet{successcolor}{accentgreen}
\colorlet{warningred}{errorcolor}
\colorlet{purple}{theoremcolor}
\colorlet{green}{accentgreen}
\colorlet{blue}{primaryblue}
\colorlet{orange}{warningcolor}
\definecolor{brown}{RGB}{121,85,72}

% Language + fonts (polyglossia)
\setdefaultlanguage{hebrew}
\setotherlanguage{english}

% Robust font selection with fallbacks (fixes Times New Roman not loadable)
% Hebrew serif
\IfFontExistsTF{David CLM}
  {\newfontfamily\hebrewfont[Script=Hebrew]{David CLM}}
  {\IfFontExistsTF{Noto Serif Hebrew}
     {\newfontfamily\hebrewfont[Script=Hebrew]{Noto Serif Hebrew}}
     {\newfontfamily\hebrewfont[Script=Hebrew]{Noto Sans Hebrew}}}

% Hebrew sans
\IfFontExistsTF{Nachlieli CLM}
  {\newfontfamily\hebrewfontsf[Script=Hebrew]{Nachlieli CLM}}
  {\IfFontExistsTF{Noto Sans Hebrew}
     {\newfontfamily\hebrewfontsf[Script=Hebrew]{Noto Sans Hebrew}}
     {\newfontfamily\hebrewfontsf[Script=Hebrew]{David CLM}}}

% Latin main font (fallbacks for missing Times New Roman)

\usepackage{unicode-math}
\setmainfont{XITS}
\setmathfont{XITSMath}


% siunitx
\sisetup{
  output-decimal-marker = {.},
  group-separator = {,},
  group-minimum-digits = 4,
  per-mode = symbol
}

% Geometry (avoid overriding in LyX header)
\geometry{
  a4paper,
  left=2cm, right=2cm, top=2.7cm, bottom=2cm,
  headheight=16pt, headsep=0.5cm, footskip=1cm
}

% Display math spacing (fix duplicates)
\AtBeginDocument{\setlength\abovedisplayskip{6pt}}
\AtBeginDocument{\setlength\abovedisplayshortskip{6pt}}
\AtBeginDocument{\setlength\belowdisplayskip{6pt}}
\AtBeginDocument{\setlength\belowdisplayshortskip{6pt}}

% Headers
\pagestyle{fancy}
\fancyhead{}
\fancyhead[R]{\nouppercase{\leftmark}}
\renewcommand{\headrulewidth}{0.5pt}

% Footnote rule (single definition)
\AtBeginDocument{
  \renewcommand\footnoterule{%
    \kern 3pt
    \hbox to \textwidth{\hfill\vrule height 0.5pt width 0.4\textwidth}
    \kern 4pt
  }%
}

% Word spacing
\spaceskip=1.3\fontdimen2\font plus 1\fontdimen3\font minus 1.5\fontdimen4\font

% Colors (deduplicated)
\definecolor{primaryblue}{RGB}{33,150,243}
\definecolor{secondaryblue}{RGB}{63,81,181}
\definecolor{accentgreen}{RGB}{76,175,80}
\definecolor{warningcolor}{RGB}{255,152,0}
\definecolor{errorcolor}{RGB}{244,67,54}
\definecolor{infocolor}{RGB}{3,169,244}
\definecolor{notecolor}{RGB}{255,193,7}
\definecolor{theoremcolor}{RGB}{156,39,176}
\definecolor{definitioncolor}{RGB}{255,87,34}
\definecolor{examplecolor}{RGB}{0,150,136}
\definecolor{lightgray}{RGB}{245,245,245}
\definecolor{darkgray}{RGB}{66,66,66}

% QED symbol
%\renewcommand{\qedsymbol}{$\blacksquare$}

% Section formatting (wrap \@ macro with makeatletter)
\makeatletter
\renewcommand*{\@seccntformat}[1]{\hspace{0.5cm}\csname the#1\endcsname\hspace{0.5cm}}
\makeatother
\titleformat{\section}{\fontsize{20}{20}\bfseries}{\thesection}{10pt}{}
\titleformat{\subsection}{\fontsize{15}{15}\bfseries}{\thesubsection}{10pt}{}
\titleformat{\subsubsection}{\bfseries}{\thesubsubsection}{10pt}{}

% RTL equation tag direction
\makeatletter
\def\maketag@@@#1{\hbox{\m@th\normalfont\LRE{#1}}}
\def\tagform@#1{\maketag@@@{(\ignorespaces#1\unskip)}}
\makeatother

% Disjoint union, large logical and
\makeatletter
\def\moverlay{\mathpalette\mov@rlay}
\def\mov@rlay#1#2{\leavevmode\vtop{%
   \baselineskip\z@skip \lineskiplimit-\maxdimen
   \ialign{\hfil$\m@th#1##$\hfil\cr#2\crcr}}}
\newcommand{\charfusion}[3][\mathord]{#1{\ifx#1\mathop\vphantom{#2}\fi\mathpalette\mov@rlay{#2\cr#3}}\ifx#1\mathop\expandafter\displaylimits\fi}
\makeatother
\newcommand{\cupdot}{\charfusion[\mathbin]{\cup}{\cdot}}
\newcommand{\bigcupdot}{\charfusion[\mathop]{\bigcup}{\cdot}}
\newcommand{\kaliLargeLand}{\mathop{\mathlarger{\mathlarger{\mathlarger{\land}}}}}

% Hebrew text helper (fix)
\newcommand{\hebtext}[1]{\texthebrew{#1}}

% Define theorem-like environments
\newtheorem{theorem}{משפט}[section]
\newtheorem{lemma}[theorem]{למה}
\newtheorem{proposition}[theorem]{הצעה}
\newtheorem{corollary}[theorem]{מסקנה}
\newtheorem{definition}[theorem]{הגדרה}
\newtheorem{example}[theorem]{דוגמה}
\newtheorem{exercise}[theorem]{תרגיל}
\newtheorem{problem}[theorem]{בעיה}
\newtheorem{remark}[theorem]{הערה}
\newtheorem{notation}[theorem]{סימון}
\newtheorem{formula}[theorem]{נוסחה}
\newtheorem{result}[theorem]{תוצאה}
\newtheorem{experiment}[theorem]{ניסוי}
\newtheorem{physicslaw}[theorem]{חוק פיזיקלי}

% Define additional physics environments
\newenvironment{physicsnote}{\begin{notebox}{}}{\end{notebox}}
\newenvironment{physicsconstant}{\begin{mainbox}{קבוע פיזיקלי}}{\end{mainbox}}
\newenvironment{physicsunits}{\begin{mainbox}{יחידות}}{\end{mainbox}}
\newenvironment{physicsapplication}{\begin{examplebox}{יישום}}{\end{examplebox}}
\newenvironment{important}{\begin{importantbox}{}}{\end{importantbox}}
\newenvironment{summary}{\begin{summarybox}{}}{\end{summarybox}}

% Define proof environment
\newenvironment{solution}{\begin{proof}[פתרון]}{\end{proof}}

% ======== TCOLORBOX CONFIGURATIONS ========
\tcbuselibrary{breakable,skins,theorems,fitting,xparse}

% Main content box
\newtcolorbox{mainbox}[2][]{
    enhanced,
    breakable,
    colback=lightgray!20,
    colframe=primaryblue,
    fonttitle=\bfseries\large,
    title={#2},
    attach boxed title to top center={yshift=-2mm},
    boxed title style={size=small, colback=primaryblue, sharp corners},
    rounded corners,
    left=6mm,
    right=6mm,
    top=4mm,
    bottom=4mm,
    drop shadow,
    #1
}

% Theorem box
\newtcolorbox{theorembox}[2][]{
    enhanced,
    breakable,
    colback=theoremcolor!10,
    colframe=theoremcolor,
    fonttitle=\bfseries\large,
    title={משפט: #2},
    attach boxed title to top left={yshift=-2mm, xshift=5mm},
    boxed title style={size=small, colback=theoremcolor, sharp corners},
    rounded corners,
    left=6mm,
    right=6mm,
    top=4mm,
    bottom=4mm,
    drop shadow,
    #1
}

% Definition box
\newtcolorbox{definitionbox}[2][]{
    enhanced,
    breakable,
    colback=definitioncolor!10,
    colframe=definitioncolor,
    fonttitle=\bfseries\large,
    title={הגדרה: #2},
    attach boxed title to top left={yshift=-2mm, xshift=5mm},
    boxed title style={size=small, colback=definitioncolor, sharp corners},
    rounded corners,
    left=6mm,
    right=6mm,
    top=4mm,
    bottom=4mm,
    drop shadow,
    #1
}

% Example box
\newtcolorbox{examplebox}[2][]{
    enhanced,
    breakable,
    colback=examplecolor!10,
    colframe=examplecolor,
    fonttitle=\bfseries\large,
    title={דוגמה: #2},
    attach boxed title to top right={yshift=-2mm, xshift=-5mm},
    boxed title style={size=small, colback=examplecolor, sharp corners},
    rounded corners,
    left=6mm,
    right=6mm,
    top=4mm,
    bottom=4mm,
    drop shadow,
    #1
}

% Problem box
\newtcolorbox{problembox}[2][]{
    enhanced,
    breakable,
    colback=primaryblue!10,
    colframe=primaryblue,
    fonttitle=\bfseries\large,
    title={בעיה: #2},
    attach boxed title to top right={yshift=-2mm, xshift=-5mm},
    boxed title style={size=small, colback=primaryblue, sharp corners},
    rounded corners,
    left=6mm,
    right=6mm,
    top=4mm,
    bottom=4mm,
    drop shadow,
    #1
}

% Solution box
\newtcolorbox{solutionbox}[1][]{
    enhanced,
    breakable,
    colback=successcolor!10,
    colframe=successcolor,
    fonttitle=\bfseries\large,
    title=פתרון,
    attach boxed title to top right={yshift=-2mm, xshift=-5mm},
    boxed title style={size=small, colback=successcolor, sharp corners},
    rounded corners,
    left=6mm,
    right=6mm,
    top=4mm,
    bottom=4mm,
    drop shadow,
    #1
}

% Note box
\newtcolorbox{notebox}[2][]{
    enhanced,
    breakable,
    colback=notecolor!15,
    colframe=notecolor,
    fonttitle=\bfseries,
    title={הערה: #2},
    attach boxed title to top left={yshift=-2mm, xshift=5mm},
    boxed title style={size=small, colback=notecolor, sharp corners},
    rounded corners,
    left=6mm,
    right=6mm,
    top=4mm,
    bottom=4mm,
    drop shadow,
    #1
}

% Warning box
\newtcolorbox{warningbox}[2][]{
    enhanced,
    breakable,
    colback=warningcolor!15,
    colframe=warningcolor,
    fonttitle=\bfseries,
    title={אזהרה: #2},
    attach boxed title to top center={yshift=-2mm},
    boxed title style={size=small, colback=warningcolor, sharp corners},
    rounded corners,
    left=6mm,
    right=6mm,
    top=4mm,
    bottom=4mm,
    drop shadow,
    #1
}

% Important box
\newtcolorbox{importantbox}[2][]{
    enhanced,
    breakable,
    colback=errorcolor!10,
    colframe=errorcolor,
    fonttitle=\bfseries,
    title={חשוב: #2},
    attach boxed title to top center={yshift=-2mm},
    boxed title style={size=small, colback=errorcolor, sharp corners},
    rounded corners,
    left=6mm,
    right=6mm,
    top=4mm,
    bottom=4mm,
    drop shadow,
    #1
}

% Summary box
\newtcolorbox{summarybox}[2][]{
    enhanced,
    breakable,
    colback=infocolor!10,
    colframe=infocolor,
    fonttitle=\bfseries\large,
    title={סיכום: #2},
    attach boxed title to top center={yshift=-2mm},
    boxed title style={size=small, colback=infocolor, sharp corners},
    rounded corners,
    left=6mm,
    right=6mm,
    top=4mm,
    bottom=4mm,
    drop shadow,
    #1
}

% Formula box
\newtcolorbox{formulabox}[2][]{
    enhanced,
    breakable,
    colback=purple!10,
    colframe=purple,
    fonttitle=\bfseries\large,
    title={נוסחה: #2},
    attach boxed title to top center={yshift=-2mm},
    boxed title style={size=small, colback=purple, sharp corners},
    rounded corners,
    left=6mm,
    right=6mm,
    top=4mm,
    bottom=4mm,
    drop shadow,
    #1
}

% Result box
\newtcolorbox{resultbox}[2][]{
    enhanced,
    breakable,
    colback=accentgreen!15,
    colframe=accentgreen,
    fonttitle=\bfseries\large,
    title={תוצאה: #2},
    attach boxed title to top center={yshift=-2mm},
    boxed title style={size=small, colback=accentgreen, sharp corners},
    rounded corners,
    left=6mm,
    right=6mm,
    top=4mm,
    bottom=4mm,
    drop shadow,
    #1
}

% Enhanced mathematical environments
\newenvironment{calculation}
  {\begin{align}}
  {\end{align}}

\newenvironment{steps}
  {\begin{enumerate}[label=שלב \arabic*:, rightmargin=1cm]}
  {\end{enumerate}}

% Custom commands for common elements
\newcommand{\highlight}[1]{\textcolor{primaryblue}{\textbf{#1}}}
% CHANGED: avoid clash with the 'important' environment
\newcommand{\imptext}[1]{\textcolor{errorcolor}{\textbf{#1}}}
\newcommand{\note}[1]{\textcolor{notecolor}{\textit{#1}}}

% Enhanced equation highlighting
\newcommand{\boxedeq}[1]{\boxed{\displaystyle #1}}
\newcommand{\redeq}[1]{\textcolor{errorcolor}{\boxed{\displaystyle #1}}}
\newcommand{\blueeq}[1]{\textcolor{primaryblue}{\boxed{\displaystyle #1}}}

\makeatother

\begin{document}
\title{סיכום - קוונטים (2)\\
 (77605)}
\maketitle
\begin{center}
כותב: ערן ריחני
\par\end{center}

\begin{center}
סמסטר א' תשפ\textquotedblright ו, האוניברסיטה העברית
\par\end{center}

\newpage{}

\tableofcontents{}
\begin{center}
\vfill{}
\par\end{center}

\begin{center}
הערות
\par\end{center}

\[
***
\]

\begin{center}
\par\end{center}

\[
***
\]

\begin{center}
\par\end{center}

\newpage{}

\section{הרצאה 4 - }

\subsection*{חזרה על שיעור קודם}

תכונות החבורה:

סגירות,יחידה,הופכי וקיבוץ.

חבורה היא רעיון אבסטרקטי וכדי לממש אותו בעולם האמיתי אנו צריכים הצגה
של החבורה.

בפרט יש את ההצגה של המרחב הוקטורי. 

בקורס שלנו נעסוק בהפעלת איברי החבורה על מרחב וקטורי ונסמן את ההצגה
כאופרטור $T\left(g\right)$ במרחב הוקטורי.

מתוך כל ההצגות ישנן את אבני הבניין של החבורות שהן החבורות הבלתי-פריקות.

מתוך החבורות הסופיות יש מספר סופי של חבורות בלתי פריקות.
משפט האורתוגונליות

\[
\sum_{g\in G}T_{ik}^{\left(\alpha\right)^{\ast}}\left(g\right)T_{jl}^{\left(\beta\right)}\left(g\right)=\frac{\left|G\right|}{\ell_{\alpha}}\delta_{\alpha\beta}\delta_{ij}\delta_{kl}
\]


\subsection{חבורות רציפות }

חבורת הסיבובים במישור $O\left(2\right)$ ניתנת לייצוג ע\textquotedbl י
המטריצה:
\[
R\left(\alpha\right)=\begin{pmatrix}\cos\left(\alpha\right) & \sin\left(\alpha\right)\\
-\sin\left(\alpha\right) & \cos\left(\alpha\right)
\end{pmatrix}
\]

ניתן גם להסתכל על המספר המרוכב $z$ ולסובב אותו ע\textquotedbl י
$e^{i\theta}$ , זה גם ייצוג של החבורה (חד מימדית).
\[
z\rightarrow ze^{i\theta}
\]

בנוסף יש עוד ייצוג במרחב הפונקציות של 
\[
f\left(\theta\right)\rightarrow f\left(\theta-\varphi\right)
\]

אלו שלוש הצגות שונות של אותה החבורה.

\subsubsection*{סיבוב קטן מאוד }

עבור סיבוב קטן מאוד: 

\[
\begin{pmatrix}\cos\left(\alpha\right) & \sin\left(\alpha\right)\\
-\sin\left(\alpha\right) & \cos\left(\alpha\right)
\end{pmatrix}\approx I+\begin{pmatrix} & \varphi\\
-\varphi
\end{pmatrix}=I+i\varphi\begin{pmatrix} & -i\varphi\\
i\varphi
\end{pmatrix}
\]

\[
e^{i\varphi}z\approx\left(1+i\varphi\right)z
\]
\[
f\left(\theta+\varphi\right)\approx\left(1+i\varphi\left(-i\frac{\partial~}{\partial\theta}\right)\right)f\left(\theta\right)
\]

באופן כללי 
\[
T\left(\varphi\right)=I+i\varphi X
\]

כאשר $X$ יכול להיות:
\[
X=\begin{pmatrix} & -i\\
i
\end{pmatrix},\quad X=1,\quad X=-i\frac{\partial~}{\partial\theta}
\]
\[
T\left(\varphi\right)=-e^{i\varphi X}
\]

$X$ היוצר של החבורה

באופן כללי: $T\left(\theta\right)=I+i\epsilon X$

\[
I=T^{\dag}T=\left(I-i\epsilon X^{\dagger}\right)\left(I+i\epsilon X\right)=I+i\epsilon\left(X-X^{\dagger}\right)+O\left(\epsilon^{2}\right)
\]

לכן $X$ הוא הרמיטי וליניארי ולכן הוא מועמד להיות אופרטור המייצג גודל
מדיד.

\[
T\left(\epsilon\right)=I+i\sum_{a}\epsilon_{a}X_{a}
\]
\[
\left[X_{a},X_{b}\right]=i\sum_{a}C_{b}X_{c}
\]


\subsubsection*{אופרטורים }

$H$ - יוצר הזזה בזמן.

$\beta$ - יוצר הזזה מרחבית.
\[
T\left(\Delta t\right)\simeq I-i\frac{\Delta t}{\hbar}H
\]

\[
\vdots
\]
\[
\mathcal{P}iH\mathcal{P^{\dagger}}=iH
\]
\[
\mathcal{P}H\mathcal{P}^{-1}=H
\]

$\mathcal{P}$ הוא ליניארי הוא לא משנה את ההמילטוניאן אך כאשר הוא
פועל על אופרטור התנע נקבל את התנע בסימן מינוס.
\[
Ti\mathcal{P}T=i\mathcal{P}
\]

$E$ לא יכול להיות שלילי (ליניארי) ולכן,

מסקנה: $\mathbb{P}$ הוא אופרטור אנטי-ליניארי\newpage{}

\section{שבוע 3 הרצאה 5 }
\date{$\date{2.11.2025}$}

\ 
משפט הווריאציה - משפט ריץ

עבור פונקציית גל אקראית $\psi$ במרחב הילברט, ערך התצפית של ההמילטוניאן
מקיים את היחס:

\[
E\equiv\frac{\left<\psi|H|\psi\right>}{\left<\psi|\psi\right>}\geqslant E_{1}
\]

כאשר $E_{1}$ היא האנרגיה של רמת היסוד.

היחס מתקיים אם ורק אם $\psi$ הוא פונקציית הגל של רמת היסוד.
הוכחה:

כיוון שהמצבים העצמיים של $H$ מהווים בסיס שלם של המרחב, ניתן לייצג
את $\psi$ כך:

\[
\psi=\sum_{k}c_{k}\phi_{k}
\]

ולכן: 
\[
\frac{\left<\psi|H-E_{1}|\psi\right>}{\left<\psi|\psi\right>}=\frac{\sum_{k=1}^{\infty}\left(E_{k}-E_{1}\right)\left|c_{k}\right|^{2}}{\sum_{k=1}^{\infty}\left|c_{k}\right|^{2}}\geqslant0
\]

כאשר השוויון האחרון נובע מסדר רמות האנרגיה מגבוה לנמוך.

$\blacksquare$

השימוש של המשפט הוא לתת חסם עליון ל-$E_{0}$.
משפט ריץ המוכלל -
\[
\left<k|\tilde{0}\right>\sim O\left(\varepsilon\right)\qquad k\neq0
\]

\[
\bar{H}-E_{0}\sim O\left(\varepsilon^{2}\right)
\]

ערך התצפית של ההמילטוניאן הוא סטציונארי (הנקודה בה הנגזרת הראשונה
מתאפסת) בסביבה של הערכים העצמיים.
הוכחה: (סאקוראי עמ' 317 )

כדי להוכיח את המשפט נשתמש בווריאציה קטנה בפונקציית הגל .

אם מוסיפים ווריאציה קטנה לפונקציית הגל 
\[
\ket{\tilde{0}}\rightarrow\ket{\tilde{0}}+\delta\ket{\tilde{0}}
\]

ונחשב את האנרגיה הממוצעת אז נקבל שגיאה מסדר גודל של $\delta^{2}$.

נניח $\psi=\phi_{n}+\delta\phi$ עם $\left<\phi|\delta\phi\right>=0$
אז במקרה הזה:
\[
E=\ldots=E_{n}+O\left(\delta\phi^{2}\right)
\]

\ \foreignlanguage{english}{The variance thorem} 
בהנתן השונות של ערך התצפית של האנרגיה
\[
\sigma^{2}\equiv\frac{\left<\psi|\left(H-E\right)^{2}|\psi\right>}{\left<\psi|\psi\right>}=\frac{\left<\psi|H^{2}|\psi\right>}{\left<\psi|\psi\right>}-E^{2}
\]

יש לפחות ערך-עצמי מדויק אחד של $H$ בתחום $\left[E-\sigma,E+\sigma\right]$.

הוכחה:

נפתח את $\psi$ באמצעות מצבים עצמיים מנורמלים של $H$. ניתן לכתוב
את השונות כך:
\[
\sigma^{2}=\sum_{i=1}^{\infty}\left(E_{i}-E\right)^{2}\left|c_{i}\right|^{2}
\]

נסמן $E_{k}$ בתור הערכים העצמיים הקרובים ביותר של $E$. אז:
\[
\sum_{i=1}^{\infty}\left(E_{i}-E\right)^{2}\left|c_{i}\right|^{2}\geqslant\left(E_{k}-E\right)^{2}\sum_{i=1}^{\infty}\left|c_{i}\right|^{2}=\left(E_{k}-E\right)^{2}
\]

ולכן $\sigma^{2}\geqslant\left(E_{k}-E\right)^{2}$. $\blacksquare$

\newpage{}

חלקיק נע בפוטנציאל כדורי $V\left(r\right)$

נסתכל על פוטנציאל מושך 
\[
V\left(r\right)=-\frac{\alpha}{r^{s}},\qquad\alpha>0
\]
ההמילטוניאן במקרה זה יהיה:
\[
\mathcal{H}=-\frac{\hbar^{2}}{2m}\nabla^{2}-\frac{\alpha}{r^{s}}
\]
 
\[
\psi\left(r\right)=Ne^{-\nicefrac{r^{2}}{r_{0}^{2}}}
\]

כאשר 
\[
N=\left(\frac{2}{\pi}\right)^{\nicefrac{3}{4}}r_{0}^{-\nicefrac{3}{2}}
\]

הוא פקטור נורמליזציה.

מתוך סימטריה נסיק:

\[
\left<x_{i}\right>=\left<p_{i}\right>=0,\quad i=1,2,3
\]

מצד שני עבור אי הוודאות נקבל:
\[
\Delta x_{i}=\sqrt{\left<x_{i}^{2}\right>-\left<x_{i}\right>^{2}}\sim r_{0}\Rightarrow\Delta p_{i}=\sqrt{\left<p_{i}^{2}\right>-\left<p_{i}\right>^{2}}=\sqrt{\left<p_{i}^{2}\right>}\sim\frac{\hbar}{r_{0}}
\]

משיקולי מימדים נכתוב
\[
\left<\psi|\frac{-\alpha}{r^{s}}|\psi\right>=\intop dr\psi^{\dag}\left(r\right)V\left(r\right)\psi\left(r\right)=
\]

נחליף משתנים $x=\frac{r}{r_{0}}$
\[
\left<\mathcal{H}\right>=\frac{\left<\boldsymbol{P}^{2}\right>}{2m}+\left<V\right>\sim\frac{\hbar^{2}}{2mr_{0}^{2}}-\frac{\alpha}{r_{0}^{2}}
\]

\end{document}
